% Lagebeziehungen – bank/lagebeziehungen/ (zusammenbau v0.4, Heft)
\einheitenkopf[t5]{Lagebeziehungen}
\setcounter{aufgabe}{48}

% Anlauf (ohne Prüfkennung)
\begin{aufgabe}{Punktprobe und Seitenlage – Anlauf}
\begin{teile}
\teil $E: 2x - y + z = 3$. Die Koordinaten von $P$ eingesetzt ergeben links $6$. Wo liegt $P$? \\ \kreuz{in $E$} \\ \kreuz{auf der Seite mit größerem Wert} \\ \kreuz{auf der Seite mit kleinerem Wert}
\teil Liegt $P(3 | 1 | 4)$ in $E: 2x - y + z = 7$?
\teil $E_1: 2x - y + 2z = 1$ und $E_2: 2x - y + 2z = 7$ sind parallel. Liegt $P(1 | 1 | 2)$ zwischen den beiden Ebenen? Begründe.
\end{teile}
\end{aufgabe}

\begin{aufgabe}{Punktprobe und Seitenlage – Prüfungsaufgaben}
\begin{teile}
\teil $E: \vec{x} = (2 | 0 | 1) + t \cdot (1 | 1 | 0) + s \cdot (0 | 1 | 2)$. Weise nach, dass $P(4 | 5 | 7)$ in $E$ liegt. \hfill (Abitur 2021 GK)
\teil $E: \vec{x} = (1 | 1 | 0) + t \cdot (2 | 0 | 1) + s \cdot (0 | 1 | 2)$. Prüfe, ob $P(5 | 3 | 5)$ in $E$ liegt. \hfill (Abitur 2021 GK)
\teil Eine Lärmschutzwand steht auf der Fahrbahn (xy-Ebene) und liegt in der Ebene $W: 5x - 2y = 0$ (1 LE = 1 m). Weise nach, dass der Punkt $K(2 | 5 | 3)$ in $W$ liegt und dass $W$ senkrecht auf der Fahrbahn steht. \hfill (Abitur 2018 GK)
\teil $P(1 | 0 | -2)$, $Q(3 | 4 | 0)$ und $E: x + 2y + z = 11$. Weise nach, dass $Q$ in $E$ liegt und dass $\vec{PQ}$ ein Normalenvektor von $E$ ist. \hfill (Abitur 2025 GK)
\teil $E: 2x + y - 2z = 3$; die Gerade $g$ geht durch $P(1 | 3 | 1)$ und $Q(5 | 5 | -3)$. Weise nach, dass $P$ in $E$ liegt, und begründe, dass $g$ senkrecht auf $E$ steht. \hfill (Abitur 2026 GK)
\teil Ein Quader hat die Kante $AE$ mit $A(5 | 0 | 0)$ und $E(5 | 0 | 3)$; die Ebene $W: x + 2y + 2z = 9$ ist gegeben. Ein Lösungsweg lautet: (1) $P(5 | 0 | r)$ mit $0 \le r \le 3$; (2) $5 + 2 \cdot 0 + 2r = 9$. Formuliere eine passende Aufgabenstellung und führe sie zu Ende. \hfill (Abitur 2024 GK)
\teil Ein Körper hat die Kante $BC$ mit $B(2 | 6 | 0)$ und $C(2 | 6 | 4)$; die Ebene $W: 3x - y + z = 1$ ist gegeben. Ein Lösungsweg lautet: (1) $P(2 | 6 | r)$ mit $0 \le r \le 4$; (2) $3 \cdot 2 - 6 + r = 1$. Formuliere eine passende Aufgabenstellung und führe sie zu Ende. \hfill (Abitur 2024 GK)
\teil Ein Quader hat die Kante $FG$ mit $F(0 | 3 | 2)$ und $G(6 | 3 | 2)$; die Ebene $W: x - y + 2z = 5$ ist gegeben. Ein Lösungsweg lautet: (1) $P(r | 3 | 2)$ mit $0 \le r \le 6$; (2) $r - 3 + 2 \cdot 2 = 5$. Formuliere eine passende Aufgabenstellung und führe sie zu Ende. \hfill (Abitur 2024 GK)
\teil Ein Körper ist der Teil des ersten Oktanten zwischen den parallelen Ebenen $L_1: x + 2y + z = 2$ und $L_2: x + 2y + z = 6$. Begründe, dass $T(1 | 0{,}5 | 1)$ im Inneren des Körpers liegt. \hfill (Abitur 2023 GK)
\teil Ein Haus ist vereinfacht der Quader mit $0 \le x \le 8$, $0 \le y \le 4$ und $0 \le z \le 5$ (1 LE = 1 m). Untersuche für die Ebenen $S: x - 2y = 0$ und $T: x + 2y = 0$, ob sie durch das Innere des Hauses verlaufen. \hfill (Abitur 2019 GK)
\end{teile}
\end{aufgabe}

% Anlauf (ohne Prüfkennung)
\begin{aufgabe}{Parameter aus der Lagebedingung – Anlauf}
\begin{teile}
\teil Für welche Zahl $k$ liegt $P(1 | 2 | 0)$ in der Ebene $E_k: kx + 3y + z = 8$? Wer gegen wen? \\ \kreuz{Punkt gegen Ebene} \\ \kreuz{Gerade gegen Ebene} \\ \kreuz{Punkt gegen Gerade}
\teil $E_k: kx + 2y - z = 5$ enthält $P(2 | 1 | 3)$. Bestimme $k$. $k = $ \leerfeld
\teil Eine Ebene der Form $r \cdot x + s \cdot z = 0$ enthält den Punkt $P(4 | 1 | -2)$. Gib passende Werte für $r$ und $s$ an. $r = $ \leerfeld , $s = $ \leerfeld
\end{teile}
\end{aufgabe}

\begin{aufgabe}{Parameter aus der Lagebedingung – Prüfungsaufgaben}
\begin{teile}
\teil Das Dreieck mit $A(3 | 0 | 0)$, $B(0 | 6 | 0)$ und $C(0 | 0 | 2)$ liegt in einer Ebene der Form $2x + y + kz = 6$. Bestimme $k$. \hfill (Abitur 2023 GK) \\ $k = $ \leerfeld
\teil Die Ebene $L$ enthält die x-Achse und den Punkt $G(3 | 2 | 5)$; ihre Gleichung hat die Form $r \cdot y + s \cdot z = 0$. Gib passende Werte für $r$ und $s$ an. \hfill (Abitur 2023 GK) \\ $r = $ \leerfeld , $s = $ \leerfeld
\end{teile}
\end{aufgabe}

% Anlauf (ohne Prüfkennung)
\begin{aufgabe}{Gerade und Ebene – Anlauf}
\begin{teile}
\teil Der allgemeine Punkt von $g$ in die Gleichung von $E$ eingesetzt ergibt $4 + 3t - 3t = 4$. Wie liegt $g$ zu $E$? \\ \kreuz{liegt in der Ebene} \\ \kreuz{echt parallel} \\ \kreuz{genau ein Schnittpunkt}
\teil Weise nach, dass $g: \vec{x} = (1 | 2 | 0) + t \cdot (1 | -1 | 1)$ in $E: x + 2y + z = 5$ liegt.
\teil $g: \vec{x} = (1 | 0 | 2) + t \cdot (2 | 1 | 1)$ und $E: x - 2y + z = 1$. Entscheide mit der Merkregel, wie $g$ zu $E$ liegt.
\end{teile}
\end{aufgabe}

\begin{aufgabe}{Gerade und Ebene – Prüfungsaufgaben}
\begin{teile}
\teil Weise nach, dass die Gerade $s: \vec{x} = (2 | -1 | 3) + r \cdot (0 | 6 | -3)$ in $E: x + y + 2z = 7$ liegt. \hfill (Abitur 2021 GK)
\teil Eine Drohne fliegt auf $h: \vec{x} = (100 | 50 | 40) + s \cdot (30 | 10 | 1)$. Die Unterseite einer Wolkendecke liegt in $E: x - 30z = -800$ (1 LE = 1 m). Weise nach, dass die Flugbahn parallel zu $E$ verläuft, und prüfe, ob sie in $E$ liegt. \hfill (Abitur 2017 GK)
\teil $P(1 | 3 | 8)$ und $Q_t(5 | t | 6)$ mit reellem $t$. Gibt es ein $t$, für das die Gerade $PQ_t$ parallel zur xy-Ebene verläuft? Begründe. \hfill (Abitur 2024 GK)
\teil $A(3 | 1 | 2)$ und $B(5 | -4 | 2)$. Begründe, dass die Gerade $AB$ parallel zur xy-Ebene verläuft. \hfill (Abitur 2017 GK)
\end{teile}
\end{aufgabe}

% Anlauf (ohne Prüfkennung)
\begin{aufgabe}{Lagebefunde im Sachzusammenhang – Anlauf}
\begin{teile}
\teil Frage: Liegt der Schatten der Lampe auf der Wand? Rechnung: Durchstoßpunkt der Lichtgeraden mit der Wandebene $(2 | 0 | 1{,}8)$. Ist die Frage damit beantwortet? \\ \kreuz{ja} \\ \kreuz{nein – ein Bereich ist noch zu prüfen}
\teil Eine Wand liegt in der Ebene $y = 6$ mit $0 \le x \le 6$ und $0 \le z \le 3$. Eine Lampe in $L(2 | 0 | 4)$ beleuchtet die Spitze $P(3 | 2 | 3)$ eines Pfostens. Liegt der Schatten von $P$ auf der Wand?
\teil Eine Terrasse ist das Rechteck mit $0 \le x \le 4$ und $0 \le y \le 6$ in der xy-Ebene. Ein Lösungsschritt lautet: $(x | y | 0) = (2 | 1 | 3) + \lambda \cdot (1 | 1 | -1)$ liefert $\lambda = 3$ und $(5 | 4 | 0)$. Erläutere den Schritt und ziehe eine Folgerung für den Schatten der Ecke $(2 | 1 | 3)$ eines Sonnensegels.
\end{teile}
\end{aufgabe}

\begin{aufgabe}{Lagebefunde im Sachzusammenhang – Prüfungsaufgaben}
\begin{teile}
\teil Sonnenlicht fällt in Richtung $(-1 | -2 | -1)$; eine Hauswand liegt in der xz-Ebene mit $0 \le x \le 6$ und $0 \le z \le 3$. Schritt I: $(x | 0 | z) = (4 | 2 | 3) + \lambda \cdot (-1 | -2 | -1)$ liefert $\lambda = 1$ und $(3 | 0 | 2)$. Schritt II: $0 < 3 < 6$ und $0 < 2 < 3$. Erläutere beide Schritte im Sachzusammenhang. \hfill (Abitur 2024 GK)
\teil Eine Wand liegt in der xz-Ebene mit $0 \le x \le 5$ und $0 \le z \le 4$; $(1 | 2 | 4)$ ist eine Ecke der Unterkante eines Rollos. Ein Lösungsschritt lautet: $(x | 0 | z) = (1 | 2 | 4) + \mu \cdot (1 | -2 | -1)$ liefert $\mu = 1$ und $(2 | 0 | 3)$. Erläutere die Bedeutung dieses Schritts und ziehe eine Folgerung. \hfill (Abitur 2023 GK)
\end{teile}
\end{aufgabe}
